% SGA 3, Expose I, Loop 1 English reconstruction.
% Scope: Section 4.4 through the end of Section 4.7, stopping immediately
% before Section 5 (Group cohomology).
% Authority: Polo--Gille Exp1-13oct24.pdf, component-local pages 23--31
% (combined-reader pages 33--41; p. 31 is included only through Section 4.7).
% Force-OCR and jcreinhold are drafting and comparison witnesses only.

% Source page: Exp1 component-local p. 23; combined-reader p. 33.
% Re-edition printed p. 23; Section 4.4 begins near the foot of the page.
\subsection*{\SGAInlineRefTarget{sga3:i:subsection:4-4}4.4. 可对角化群}

\medskip
\noindent\SGAInlineRefTarget{sga3:i:paragraph:4-4-1}\textbf{4.4.1.}
\(\mathbb G_m\) 的构造可以推广如下. 设 \(M\) 是一个交换群,
\(M_{\mathbb Z}\) 是按照 \SGARefLink{sga3:i:subsection:4-1}{4.1}
的方法与之关联的群概形. 考虑下式定义的函子:
\[
  \mathbf D(M)(S)
    =\operatorname{Hom}_{\mathrm{groups}}(M,\mathbb G_m(S))
    \simeq
      \underline{\operatorname{Hom}}_{S\text{-}\mathbf{Gr}}
        (M_S,\mathbb G_{m,S}).
\]

% Source page transition: Exp1 component-local p. 24;
% combined-reader p. 34; re-edition printed p. 24.
这是一个交换 \(\widehat{(\mathbf{Sch})}\)-群, 由一个记为 \(D(M)\)
的群概形表示; 因而按照定义,
\[
  D(M)\simeq
  \underline{\operatorname{Hom}}_{(\mathbf{Sch})\text{-}\mathbf{Gr}}
    (M_{\mathbb Z},\mathbb G_m).
\]
事实上, 令
\[
  D(M)=\operatorname{Spec}\mathbb Z[M],
\]
其中 \(\mathbb Z[M]\) 是 \(M\) 在 \(\mathbb Z\) 上的群代数.
由代数 \(\mathbb Z[M]\) 的定义本身可知
\[
  D(M)(S)
    =\operatorname{Hom}_{\mathrm{Alg.}}
       \bigl(\mathbb Z[M],\Gamma(S,\mathcal O_S)\bigr)
    \simeq
      \operatorname{Hom}_{\mathrm{groups}}
       \bigl(M,\Gamma(S,\mathcal O_S)^{\times}\bigr)
\]
成立.

\medskip
\noindent\SGAInlineRefTarget{sga3:i:paragraph:4-4-2}\textbf{4.4.2.}
因此, 对每个概形 \(S\), 都有一个在 \(S\) 上仿射的 \(S\)-群概形,
\[
  D_S(M)=D(M)_S
    =\bigl(
      \underline{\operatorname{Hom}}_{(\mathbf{Sch})\text{-}\mathbf{Gr}}
        (M_{\mathbb Z},\mathbb G_m)
      \bigr)_S
    =\underline{\operatorname{Hom}}_{S\text{-}\mathbf{Gr}}
       (M_S,\mathbb G_{m,S}).
\]
它对应于 \(\mathcal O_S\)-双代数 \(\mathcal O_S[M]\),
其对角映射和增广由下式定义:
\[
  \Delta(x)=x\otimes x,
  \qquad
  \varepsilon(x)=1,
  \qquad x\in M.
\]

\medskip
\noindent\SGAInlineRefTarget{sga3:i:paragraph:4-4-3}\textbf{4.4.3.}
若 \(f:M\to N\) 是交换群同态, 则可用显然的方式定义 \(S\)-群态射
\[
  D_S(f):D_S(N)\longrightarrow D_S(M),
\]
从而得到一个从交换群范畴到在 \(S\) 上仿射的 \(S\)-群范畴的函子
\[
  D_S:M\longmapsto D_S(M)
\]
它也可定义为函子 \(M\mapsto M_S\) 与函子
\[
  M_S\longmapsto
  \underline{\operatorname{Hom}}_{S\text{-}\mathbf{Gr}}
    (M_S,\mathbb G_{m,S}).
\]
的复合. 这个函子与基扩张可交换.

与形如 \(D_S(M)\) 的群同构的 \(S\)-群称为\emph{可对角化群}.
注意, \(M\) 的元素可解释为 \(D_S(M)\) 的某些特征标,
即下列集合的某些元素:
\[
  \operatorname{Hom}_{S\text{-}\mathbf{Gr}}
    (D_S(M),\mathbb G_{m,S}).
\]
(比较 VIII, 1.)

\medskip
\noindent\SGAInlineRefTarget{sga3:i:paragraph:4-4-4}\textbf{4.4.4.}
下面给出一些可对角化群的例子. 首先,
\[
  D(\mathbb Z)=\mathbb G_m,
  \qquad
  D(\mathbb Z^r)=(\mathbb G_m)^r.
\]
令
\[
  \boldsymbol\mu_n=D(\mathbb Z/n\mathbb Z),
\]
并称之为 \emph{\(n\) 次单位根群}. 事实上,
\[
  \boldsymbol\mu_n(S)
    =\operatorname{Hom}_{\mathrm{groups}}
       \bigl(\mathbb Z/n\mathbb Z,\Gamma(S,\mathcal O_S)^{\times}\bigr)
    =\{f\in\Gamma(S,\mathcal O_S)\mid f^n=1\}.
\]
\(S\)-群 \(\boldsymbol\mu_{n,S}\) 对应于 \(\mathcal O_S\)-代数
\(\mathcal O_S[T]/(T^n-1)\). 特别地, 假设 \(S\) 是一个特征为
\(p=n\) 的域 \(k\) 的谱. 令 \(T-1=s\), 得到
\[
  k[T]/(T^p-1)=k[s]/(s^p),
\]
这表明 \(\boldsymbol\mu_{p,k}\) 的底拓扑空间只有一个点,
其局部环为 Artin \(k\)-代数 \(k[s]/(s^p)\). 同样地, 注意
\(S\)-概形 \(\mathbb G_{a,S}\), \(\mathbb G_{m,S}\) 和
\(\mathbf O_S\) 都在 \(S\) 上光滑, 而 \(D_S(M)\) 在

% Source page transition: Exp1 component-local p. 25;
% combined-reader p. 35; re-edition printed p. 25.
\(S\) 上平坦. 此外, 它在 \(S\) 上形式光滑 (分别地, 光滑)
当且仅当 \(S\) 的任何剩余特征都不整除 \(M\) 的挠元的阶
(分别地, 并且还要求 \(M\) 为有限型); 参见
\SGARefLink{sga3:viii:proposition:2:1}{第 VIII 篇, 2.1}.

\subsection*{\SGAInlineRefTarget{sga3:i:subsection:4-5}4.5. 群的其他例子}

前述方法适用于``典型群'' (线性群 \(\operatorname{GL}_n\), 辛群
\(\operatorname{Sp}_n\), 正交群 \(\operatorname O_n\) 等).
例如, 将 \(\operatorname{GL}_n\) 定义为表示下述函子
\(\mathbf{GL}_n\) 的群概形:
\[
  \mathbf{GL}_n(S)
    =\operatorname{GL}\bigl(n,\Gamma(S,\mathcal O_S)\bigr)
    =\operatorname{Aut}_{\mathcal O_S}(\mathcal O_S^n).
\]
例如, 可将其构造为 \(\operatorname{Spec}\mathbb Z[T_{ij}]\),
\(1\leq i,j\leq n\), 中由函数
\[
  f=\det\bigl((T_{ij})_{i,j=1}^n\bigr),
\]
定义的开子概形, 即 \(\operatorname{Spec}\mathbb Z[T_{ij},f^{-1}]\).

\subsection*{\SGAInlineRefTarget{sga3:i:subsection:4-6}4.6. 概形范畴中的函子模}

对概形 \(S\) 上的每个 \(\mathcal O_S\)-模, 我们都要关联一个
\(\mathbf O_S\)-模, 其中 \(\mathbf O_S\) 表示在 4.3.3 中引入的函子环.
这有以下两种不同的做法.

\medskip
\noindent\SGAInlineRefTarget{sga3:i:definition:4-6-1}\textbf{定义 4.6.1.}
设 \(S\) 是一个概形. 对每个 \(\mathcal O_S\)-模 \(\mathscr F\),
用 \(\mathbf V(\mathscr F)\) 和 \(\mathbf W(\mathscr F)\) 表示
\((\mathbf{Sch})_{/S}\) 上由下式定义的反变函子:
\begin{align*}
  \mathbf V(\mathscr F)(S')
    &=\operatorname{Hom}_{\mathcal O_{S'}}
       (\mathscr F\otimes_{\mathcal O_S}\mathcal O_{S'},
        \mathcal O_{S'}),\\
  \mathbf W(\mathscr F)(S')
    &=\Gamma(S',\mathscr F\otimes_{\mathcal O_S}\mathcal O_{S'}),
\end{align*}
其中 \(\mathscr F\otimes_{\mathcal O_S}\mathcal O_{S'}\) 表示
\(\mathcal O_S\)-模 \(\mathscr F\) 在 \(S'\) 上的逆像.

于是 \(\mathbf V(\mathscr F)\) 和 \(\mathbf W(\mathscr F)\)
以显然的方式带有 \(\mathbf O_S\)-模结构 (回顾
\(\mathbf O_S(S')=\Gamma(S',\mathcal O_{S'})
=\mathbf W(\mathcal O_S)(S')\)). 因而实际上已经定义了两个
从 \(\mathcal O_S\)-模范畴到 \(\mathbf O_S\)-模范畴的函子
\(\mathbf V\) 和 \(\mathbf W\), 其中 \(\mathbf V\) 反变,
\(\mathbf W\) 协变.

在本段余下部分中, 我们限于所涉及的 \(\mathcal O_S\)-模都是拟凝聚的情形.
因此将 \(\mathbf V\) 和 \(\mathbf W\) 视为以下函子:
\[
  \mathbf V:(\mathcal O_S\text{-}\mathbf{Mod.q.c.})^{\circ}
    \longrightarrow(\mathbf O_S\text{-}\mathbf{Mod.}),
  \qquad
  \mathbf W:(\mathcal O_S\text{-}\mathbf{Mod.q.c.})
    \longrightarrow(\mathbf O_S\text{-}\mathbf{Mod.}).
\]

\medskip
\noindent\SGAInlineRefTarget{sga3:i:remark:4-6-1-1}\textbf{注 4.6.1.1.}
\footnote{编者注: 为便于后文引用, 添加了编号
\SGARefLink{sga3:i:remark:4-6-1-1}{4.6.1.1}.}
读者会注意到, 在下文中, 仅涉及函子 \(\mathbf W\) 的每个命题
对任意 \(\mathcal O_S\)-模仍成立, 而不必要求其拟凝聚.

\medskip
\noindent\SGAInlineRefTarget{sga3:i:proposition:4-6-2}\textbf{命题 4.6.2.}
\begin{enumerate}
  \item[\SGAInlineRefTarget{sga3:i:proposition:4-6-2:item:i}(i)]
  函子 \(\mathbf V\) 和 \(\mathbf W\) 与基扩张可交换:
  若 \(S'\) 在 \(S\) 上, 且 \(\mathscr F\) 是拟凝聚
  \(\mathcal O_S\)-模, 则
  \[
    \mathbf V(\mathscr F\otimes\mathcal O_{S'})
      \simeq\mathbf V(\mathscr F)_{S'},
    \qquad
    \mathbf W(\mathscr F\otimes\mathcal O_{S'})
      \simeq\mathbf W(\mathscr F)_{S'}.
  \]

% Source page transition: Exp1 component-local p. 26;
% combined-reader p. 36; re-edition printed p. 26.
  \item[\SGAInlineRefTarget{sga3:i:proposition:4-6-2:item:ii}(ii)]
  函子 \(\mathbf V\) 和 \(\mathbf W\) 是全忠实的: 典范映射
  \[
    \operatorname{Hom}_{\mathcal O_S}(\mathscr F,\mathscr F')
      \longrightarrow
    \operatorname{Hom}_{\mathbf O_S}
      (\mathbf V(\mathscr F'),\mathbf V(\mathscr F)),
  \]
  \[
    \operatorname{Hom}_{\mathcal O_S}(\mathscr F,\mathscr F')
      \longrightarrow
    \operatorname{Hom}_{\mathbf O_S}
      (\mathbf W(\mathscr F),\mathbf W(\mathscr F'))
  \]
  都是双射.

  \item[\SGAInlineRefTarget{sga3:i:proposition:4-6-2:item:iii}(iii)]
  函子 \(\mathbf V\) 和 \(\mathbf W\) 是加性的:
  \[
    \mathbf V(\mathscr F\oplus\mathscr F')
      \simeq
    \mathbf V(\mathscr F)\times_S\mathbf V(\mathscr F'),
    \qquad
    \mathbf W(\mathscr F\oplus\mathscr F')
      \simeq
    \mathbf W(\mathscr F)\times_S\mathbf W(\mathscr F').
  \]
\end{enumerate}

\SGARefLink{sga3:i:proposition:4-6-2:item:i}{(i)} 与
\SGARefLink{sga3:i:proposition:4-6-2:item:iii}{(iii)} 由定义立即可得.
对于 \SGARefLink{sga3:i:proposition:4-6-2:item:ii}{(ii)}, 取 \(S'\)
为 \(S\) 的开子概形即可. 证明留给读者
(对于 \(\mathbf V\), 使用 EGA~II, 1.7.14).

回顾一下 (参见 \SGARefLink{sga3:i:definition:3-1-4}{3.1.4}), 若
\(\mathbf F\) 和 \(\mathbf F'\) 是 \(\mathbf O_S\)-模, 则
\(\underline{\operatorname{Hom}}_{\mathbf O_S}(\mathbf F,\mathbf F')\)
表示取值于交换群的 \(S\)-函子, 它对每个 \(S'\to S\) 赋予群
\[
  \operatorname{Hom}_{\mathbf O_{S'}}(\mathbf F_{S'},\mathbf F'_{S'}).
\]

\medskip
\noindent\SGAInlineRefTarget{sga3:i:proposition:4-6-3}\textbf{命题 4.6.3.}
\textit{在 \((\mathbf O_S\text{-}\mathbf{Mod.})\) 中有以下典范态射:}
\footnote{编者注: 新增了同构
\(\underline{\operatorname{Hom}}_{\mathbf O_S}
(\mathbf W(\mathscr F),\mathbf W(\mathscr F'))
\simeq
\underline{\operatorname{Hom}}_{\mathbf O_S}
(\mathbf V(\mathscr F'),\mathbf V(\mathscr F))\).}

\begin{figure}[H]
  \centering
  \begin{tikzcd}[column sep=3.2em,row sep=3.5em]
  \underline{\operatorname{Hom}}_{\mathbf O_S}
    \bigl(\mathbf W(\mathscr F),\mathbf W(\mathscr F')\bigr)
    \arrow[rr,equal]
  &&
  \underline{\operatorname{Hom}}_{\mathbf O_S}
    \bigl(\mathbf V(\mathscr F'),\mathbf V(\mathscr F)\bigr)
  \\
  &
  \mathbf W\!\left(
    \mathscr Hom_{\mathcal O_S}(\mathscr F,\mathscr F')
  \right)
    \arrow[ul]
    \arrow[ur]
  &
  {} .
\end{tikzcd}

  \phantomsection\label{fig:sga3-exp1-p026-prop463-canonical-morphisms}
\end{figure}

这由 \SGARefLink{sga3:i:proposition:4-6-2:item:i}{4.6.2(i)} 和
\SGARefLink{sga3:i:proposition:4-6-2:item:ii}{(ii)} 立即可得.

\medskip
\noindent\SGAInlineRefTarget{sga3:i:notation:4-6-3-1}\textbf{记号 4.6.3.1.}
设 \(\mathscr F\) 是拟凝聚 \(\mathcal O_S\)-模. 根据 EGA~II, 1.7.8,
\(S\)-函子 \(\mathbf V(\mathscr F)\) 可由一个在 \(S\) 上仿射的
\(S\)-概形表示, 记为 \(\mathbb V(\mathscr F)\), 并称为由
\(\mathscr F\) 定义的\emph{向量纤维化 (vector fibration)}:
\footnote{编者注: 已将``向量丛''替换为``向量纤维化''.
按照现行用法, 秩为 \(r\) 的局部平凡向量纤维化称为
``秩为 \(r\) 的向量丛'', 即其截面层局部同构于
\(\mathcal O_S^{\oplus r}\).}
\[
  \mathbb V(\mathscr F)=\operatorname{Spec}(\mathcal S(\mathscr F)),
\]
其中 \(\mathcal S(\mathscr F)\) 表示 \(\mathcal O_S\)-模
\(\mathscr F\) 的对称代数.
\footnote{编者注: 文章 [Ni04] 和 [Ni02] 分别表明, 若 \(S\)
是 Noether 概形, \(\mathscr F\) 是凝聚 \(\mathcal O_S\)-模,
则 \(\mathbf W(\mathscr F)\), 分别地, 对每个 \(T\to S\)
赋予 \(\operatorname{Aut}_{\mathcal O_T}
(\mathscr F\otimes\mathcal O_T)\) 的 \(S\)-群, 可表当且仅当
\(\mathscr F\) 局部自由.}

\medskip
\noindent\SGAInlineRefTarget{sga3:i:proposition:4-6-4}\textbf{命题 4.6.4.}
\textit{设 \(\mathscr F\) 与 \(\mathscr F'\) 是两个拟凝聚
\(\mathcal O_S\)-模, \(\mathscr A\) 是一个拟凝聚 \(\mathcal O_S\)-代数.
有一个具有函子性的同构
\[
  \operatorname{Hom}_S
   \bigl(\operatorname{Spec}(\mathscr A),
     \underline{\operatorname{Hom}}_{\mathbf O_S}
       (\mathbf W(\mathscr F'),\mathbf W(\mathscr F))\bigr)
  \xrightarrow{\ \sim\ }
  \operatorname{Hom}_{\mathcal O_S}
    (\mathscr F',\mathscr F\otimes_{\mathcal O_S}\mathscr A).
\]}

% Source page transition: Exp1 component-local p. 27;
% combined-reader p. 37; re-edition printed p. 27.
事实上, 令 \(X=\operatorname{Spec}(\mathscr A)\). 左端典范同构于
\[
  \underline{\operatorname{Hom}}_{\mathbf O_S}
    (\mathbf W(\mathscr F'),\mathbf W(\mathscr F))(X),
\]
即按照定义, 典范同构于
\begin{align*}
  &\operatorname{Hom}_{\mathbf O_X}
    (\mathbf W(\mathscr F')_X,\mathbf W(\mathscr F)_X)\\
  &\qquad\simeq
  \operatorname{Hom}_{\mathbf O_X}
    \bigl(\mathbf W(\mathscr F'\otimes\mathcal O_X),
          \mathbf W(\mathscr F\otimes\mathcal O_X)\bigr)
\end{align*}
(参见 \SGARefLink{sga3:i:proposition:4-6-2:item:i}{4.6.2(i)}).
由 \SGARefLink{sga3:i:proposition:4-6-2:item:ii}{4.6.2(ii)}, 这也可写为
\[
  \operatorname{Hom}_{\mathcal O_X}
    (\mathscr F'\otimes\mathcal O_X,\mathscr F\otimes\mathcal O_X)
  =
  \operatorname{Hom}_{\mathcal O_S}
    (\mathscr F',\pi_*\pi^*(\mathscr F)),
\]
其中 \(\pi:X\to S\) 是结构态射. 而由 EGA~II, 1.4.7, 有
\[
  \pi_*\pi^*(\mathscr F)\simeq\mathscr F\otimes\mathscr A,
\]
证明完毕.

\medskip
\noindent\SGAInlineRefTarget{sga3:i:corollary:4-6-4-1}\textbf{推论 4.6.4.1.}
\textit{有一个典范同构
\[
  \mathbf W(\mathscr F\otimes\mathscr A)
    \simeq
  \underline{\operatorname{Hom}}_S
    (\operatorname{Spec}(\mathscr A),\mathbf W(\mathscr F)).
\]}

事实上,
\footnote{编者注: 以下对原文作了展开.}
设 \(f:S'\to S\) 是一个 \(S\)-概形, 并令 \(X'=X\times_S S'\).
有一个 笛卡尔方块

\begin{figure}[H]
  \centering
  \begin{tikzcd}[column sep=4.0em,row sep=3.0em]
  X'
    \arrow[r,"f'"]
    \arrow[d,"\pi'"']
  &
  X
    \arrow[d,"\pi"]
  \\
  S'
    \arrow[r,"f"]
  &
  S
\end{tikzcd}

  \phantomsection\label{fig:sga3-exp1-p027-cor4641-cartesian-square}
\end{figure}

由 EGA~II, 1.5.2, \(X'\) 在 \(S'\) 上仿射, 并且
\(\pi'_*\mathcal O_{X'}=f^*(\mathscr A)\). 因此
\begin{align*}
  &\underline{\operatorname{Hom}}_S
     (\operatorname{Spec}(\mathscr A),\mathbf W(\mathscr F))(S')\\
  &\qquad=
  \operatorname{Hom}_{S'}
     \bigl(\operatorname{Spec}(f^*\mathscr A),
           \mathbf W(f^*\mathscr F)\bigr).
\end{align*}
将 \SGARefLink{sga3:i:proposition:4-6-4}{4.6.4} 应用于
\(f^*\mathscr F\), \(\mathscr F'=\mathcal O_{S'}\) 和
\(f^*\mathscr A\), 得知上式等于
\begin{align*}
  \Gamma(S',f^*\mathscr F\otimes f^*\mathscr A)
    &=\Gamma(S',f^*(\mathscr F\otimes\mathscr A))\\
    &=\mathbf W(\mathscr F\otimes\mathscr A)(S').
\end{align*}

\medskip
\noindent\SGAInlineRefTarget{sga3:i:proposition:4-6-5}\textbf{命题 4.6.5.}
\textit{若 \(\mathscr F\) 与 \(\mathscr F'\) 是两个有限型局部自由
\(\mathcal O_S\)-模, 则 \SGARefLink{sga3:i:proposition:4-6-3}{4.6.3}
中的态射都是同构.}

事实上, 对每个 \(S'\to S\),
\begin{align*}
  &\mathbf W
    \bigl(\mathscr{H}om_{\mathcal O_S}(\mathscr F,\mathscr F')\bigr)(S')\\
  &\quad=
    \Gamma\bigl(S',
      \mathscr{H}om_{\mathcal O_S}(\mathscr F,\mathscr F')
        \otimes\mathcal O_{S'}\bigr)\\
  &\quad=
    \operatorname{Hom}_{\mathcal O_{S'}}
      (\mathscr F\otimes\mathcal O_{S'},
       \mathscr F'\otimes\mathcal O_{S'}).
\end{align*}
根据 \SGARefLink{sga3:i:proposition:4-6-2:item:i}{4.6.2(i)} 和
\SGARefLink{sga3:i:proposition:4-6-2:item:ii}{(ii)}, 最后一项既同构于
\[
  \underline{\operatorname{Hom}}_{\mathbf O_S}
    (\mathbf W(\mathscr F),\mathbf W(\mathscr F'))(S')
\]
也同构于
\[
  \underline{\operatorname{Hom}}_{\mathbf O_S}
    (\mathbf V(\mathscr F'),\mathbf V(\mathscr F))(S')
\]

\medskip
\noindent\SGAInlineRefTarget{sga3:i:corollary:4-6-5-1}\textbf{推论 4.6.5.1.}
\textit{设 \(\mathscr F\) 是有限型局部自由 \(\mathcal O_S\)-模. 令
\(\mathscr F^{\vee}
=\mathscr{H}om_{\mathcal O_S}(\mathscr F,\mathcal O_S)\).
有以下典范同构
\[
  \mathbf W(\mathscr F^{\vee})
    \simeq
  \underline{\operatorname{Hom}}_{\mathbf O_S}
    (\mathbf W(\mathscr F),\mathbf O_S)
    \simeq\mathbf V(\mathscr F),
\]
\[
  \mathbf V(\mathscr F^{\vee})
    \simeq
  \underline{\operatorname{Hom}}_{\mathbf O_S}
    (\mathbf V(\mathscr F),\mathbf O_S)
    \simeq\mathbf W(\mathscr F).
\]}

最后, 我们有如下命题.

\medskip
\noindent\SGAInlineRefTarget{sga3:i:proposition:4-6-6}\textbf{命题 4.6.6.}
\textit{设 \(f:\mathscr F\to\mathscr F'\) 是有限秩局部自由
\(\mathcal O_S\)-模的态射. 态射
\[
  \mathbf W(f):\mathbf W(\mathscr F)\longrightarrow\mathbf W(\mathscr F')
\]
为单态射当且仅当 \(f\) 局部地将 \(\mathscr F\) 等同于
\(\mathscr F'\) 的一个直和因子.}

% Source page transition: Exp1 component-local p. 28;
% combined-reader p. 38; re-edition printed p. 28.
正向断言实质上已包含于 EGA~0\textsubscript{I}, 5.5.5 中.
\footnote{编者注: 下面一句作了展开.}
反之, 若 \(\mathscr F\) 局部为 \(\mathscr F'\) 的直和因子,
则对每个 \(\pi:S'\to S\), 模 \(\pi^*\mathscr F\) 都是
\(\pi^*\mathscr F'\) 的子模, 因为局部作为直和因子时结论成立.
因此
\[
  \mathbf W(\mathscr F)(S')=\Gamma(S',\pi^*\mathscr F)
\]
是
\[
  \mathbf W(\mathscr F')(S')=\Gamma(S',\pi^*\mathscr F').
\]
的子模.

\subsection*{\SGAInlineRefTarget{sga3:i:subsection:4-7}4.7. \(G\)-\(\mathcal O_S\)-模范畴}

设 \(G\) 是一个 \(S\)-群, \(\mathscr F\) 是一个
\(\mathcal O_S\)-模; \(\mathbf W(\mathscr F)\) 带有
\(\mathbf O_S\)-模结构.

\medskip
\noindent\SGAInlineRefTarget{sga3:i:definition:4-7-1}\textbf{定义 4.7.1.}
\(\mathscr F\) 上的一个 \emph{\(G\)-\(\mathcal O_S\)-模结构},
是指 \(\mathbf W(\mathscr F)\) 上的一个
\(\mathbf h_G\)-\(\mathbf O_S\)-模结构
(参见 \SGARefLink{sga3:i:definition:3-2}{3.2}).
\(G\)-\(\mathcal O_S\)-模的态射按定义是相应的
\(\mathbf h_G\)-\(\mathbf O_S\)-模的态射.
由此得到范畴 \((G\text{-}\mathcal O_S\text{-}\mathbf{Mod.})\);
我们用 \((G\text{-}\mathcal O_S\text{-}\mathbf{Mod.q.c.})\)
表示其由作为 \(\mathcal O_S\)-模是拟凝聚的那些
\(G\)-\(\mathcal O_S\)-模组成的全子范畴.

由 \SGARefLink{sga3:i:definition:3-2}{3.2}, 在 \(\mathscr F\) 上
给定 \(G\)-\(\mathcal O_S\)-模结构因而等价于给定一个
\(\widehat{(\mathbf{Sch})}_{/S}\)-群态射
\[
  \rho:\mathbf h_G\longrightarrow
    \underline{\operatorname{Aut}}_{\mathbf O_S}
      (\mathbf W(\mathscr F)).
\]

\medskip
\noindent\SGAInlineRefTarget{sga3:i:remark:4-7-1-0}\textbf{注 4.7.1.0.}
\footnote{编者注: 此注由 [DG], II,
\SGARefLink{sga3:ii:section:2}{第 2 节},
\SGARefLink{sga3:i:subsection:1-1}{1.1} 补入.}
由于 \SGARefLink{sga3:i:proposition:4-6-2}{4.6.2} 给出取值于群的
\(S\)-函子之间的一个反同构
\[
  (\dagger)\qquad
  \underline{\operatorname{Aut}}_{\mathbf O_S}
    (\mathbf W(\mathscr F))
  \simeq
  \underline{\operatorname{Aut}}_{\mathbf O_S}
    (\mathbf V(\mathscr F)),
\]
所以在 \(\mathbf W(\mathscr F)\) 或 \(\mathbf V(\mathscr F)\)
上给定 \(\mathbf h_G\)-\(\mathbf O_S\)-模结构是等价的.
事实上, 设
\[
  \rho:\mathbf h_G\longrightarrow
  \underline{\operatorname{Aut}}_{\mathbf O_S}
    (\mathbf W(\mathscr F))
\]
如上所述. 对每个 \(T\to S\) 和 \(g\in G(T)\),
用 \(\rho^*(g)\) 表示 \(\rho(g)\) 在反同构 \((\dagger)\)
下的像. 则
\[
  \rho^*(gh)=\rho^*(h)\circ\rho^*(g),
\]
因此 \(\rho^*\) 定义了 \(\mathbf V(\mathscr F)\) 上的一个
\emph{右} \(\mathbf h_G\)-\(\mathbf O_S\)-模结构. 令
\[
  \rho^{\vee}(g)=\rho^*(g^{-1}),
\]
就得到 \(\mathbf V(\mathscr F)\) 上的一个
\(\mathbf h_G\)-\(\mathbf O_S\)-模结构, 给定这一结构等价于给定
\(\rho\).

\medskip
\noindent\SGAInlineRefTarget{sga3:i:remark:4-7-1-1}\textbf{注 4.7.1.1.}
可以说, 刚才定义的范畴是由下面的 笛卡尔方块构造出来的.

\begin{figure}[H]
  \centering
  \begin{tikzcd}[column sep=3.3em,row sep=3.2em]
  (G\text{-}\mathcal O_S\text{-}\mathbf{Mod.q.c.})
    \arrow[r,hook]
    \arrow[d]
  &
  (G\text{-}\mathcal O_S\text{-}\mathbf{Mod.})
    \arrow[r]
    \arrow[d]
  &
  (\mathbf h_G\text{-}\mathbf O_S\text{-}\mathbf{Mod.})
    \arrow[d,"\text{遗忘}"]
  \\
  (\mathcal O_S\text{-}\mathbf{Mod.q.c.})
    \arrow[r,hook]
  &
  (\mathcal O_S\text{-}\mathbf{Mod.})
    \arrow[r,"\mathbf W"]
  &
  (\mathbf O_S\text{-}\mathbf{Mod.})
  \qquad .
\end{tikzcd}

  \phantomsection\label{fig:sga3-exp1-p028-rem4711-category-squares}
\end{figure}

范畴 \((\mathcal O_S\text{-}\mathbf{Mod.})\) 和
\((\mathbf O_S\text{-}\mathbf{Mod.})\) 是阿贝尔范畴, 但须记住,
一般来说函子 \(\mathbf W\) 既非左正合, 也非右正合.
\footnote{编者注: 原文中关于范畴
\((G\text{-}\mathcal O_S\text{-}\mathbf{Mod.})\) 是阿贝尔范畴的断言
已被删去, 以作更正; 参见下面的
\SGARefLink{sga3:i:corollary:4-7-2-1}{4.7.2.1}.}

% Source page transition: Exp1 component-local p. 29;
% combined-reader p. 39; re-edition printed p. 29.
\medskip
\noindent\SGAInlineRefTarget{sga3:i:remark:4-7-1-2}\textbf{注 4.7.1.2.}
\footnote{编者注: 此注为新增内容.}
设 \(\mathscr F\) 是一个 \(G\)-\(\mathcal O_S\)-模.
其不变量子层 \(\mathscr F^G\) 定义如下: 对 \(S\) 的每个开子集
\(U\),
\[
  \mathscr F^G(U)=\mathbf W(\mathscr F)^G(U)
   =\bigl\{x\in\mathscr F(U)\mid
       g\cdot x_{S'}=x_{S'}
       \text{ 对每个 }S'\xrightarrow{f}U,
       \ g\in G(S')\bigr\},
\]
其中 \(x_{S'}\) 表示 \(x\) 在
\[
  \Gamma(S',f^*\mathscr F)=\Gamma(U,f_*f^*\mathscr F).
\]
中的像.
应当注意, 自然态射
\[
  \mathbf W(\mathscr F^G)\longrightarrow\mathbf W(\mathscr F)^G
\]
一般不是同构. 例如, 令 \(S=\operatorname{Spec}(\mathbb Z)\),
并令 \(G\) 是常群 \(\mathbb Z/2\mathbb Z=\{1,\tau\}\), 在
\(\mathscr F=\mathcal O_S\) 上通过 \(\tau\cdot1=-1\) 作用.
则 \(\mathscr F^G=0\), 但若 \(R\) 是一个 \(\mathbb F_2\)-代数, 则
\[
  \mathbf W(\mathscr F)^G(\operatorname{Spec}(R))=R.
\]

\medskip
\noindent\SGAInlineRefTarget{sga3:i:paragraph:4-7-2}\textbf{4.7.2.}
从现在起直到 \SGARefLink{sga3:i:subsection:4-7}{4.7} 结束,
假设 \(G\) 在 \(S\) 上仿射.
\footnote{编者注: 将 \SGARefLink{sga3:i:proposition:4-7-2}{4.7.2}
的结果推广到 \(G\) 未必仿射, 但假设 \(G\) 与 \(\mathscr F\)
在 \(S\) 上平坦的情形, 参见 VI\textsubscript{B},
\SGARefLink{lem:sga3-VIB-11.1}{11.1}--\SGARefLink{par:sga3-VIB-11.6}{11.6} 节.}
由 \SGARefLink{sga3:i:proposition:4-6-4}{4.6.4}, 给定 \(S\)-函子态射
\[
  \rho:\mathbf h_G\longrightarrow
    \underline{\operatorname{End}}_{\mathbf O_S}
      (\mathbf W(\mathscr F))
\]
等价于给定一个 \(\mathcal O_S\)-模态射
\[
  \mu:\mathscr F\longrightarrow
    \mathscr F\otimes_{\mathcal O_S}\mathscr A(G).
\]
此时, \(\rho\) 为 \(\widehat{(\mathbf{Sch})}_{/S}\)-群态射
等价于 \(\mu\) 满足下述公理.

\medskip
\noindent\textbf{(CM 1)} 下图交换.

\begin{figure}[H]
  \centering
  \begin{tikzcd}[column sep=5.3em,row sep=3.5em]
  \mathscr F
    \arrow[r,"\mu"]
    \arrow[d,"\mu"']
  &
  \mathscr F\otimes_{\mathcal O_S}\mathscr A(G)
    \arrow[d,"\operatorname{Id}\otimes\Delta"]
  \\
  \mathscr F\otimes_{\mathcal O_S}\mathscr A(G)
    \arrow[r,"\mu\otimes\operatorname{Id}"]
  &
  \mathscr F\otimes_{\mathcal O_S}\mathscr A(G)
    \otimes_{\mathcal O_S}\mathscr A(G)
  \qquad .
\end{tikzcd}

  \phantomsection\label{fig:sga3-exp1-p029-cm1-coaction-square}
\end{figure}

\noindent\textbf{(CM 2)} 下面的复合是恒等映射.

\begin{figure}[H]
  \centering
  \begin{tikzcd}[column sep=4.5em]
  \mathscr F
    \arrow[r,"\mu"]
  &
  \mathscr F\otimes\mathscr A(G)
    \arrow[r,"\operatorname{Id}\otimes\varepsilon"]
  &
  \mathscr F\otimes\mathcal O_S
    \arrow[r,"\sim"]
  &
  \mathscr F
  \qquad .
\end{tikzcd}

  \phantomsection\label{fig:sga3-exp1-p029-cm2-counit-composite}
\end{figure}

这些公理 (CM 1) 和 (CM 2) 就是双代数 \(\mathscr A(G)\) 上的
右\emph{余模 (comodule)} 结构的公理.
\footnote{编者注: 左 \(G\)-\(\mathcal O_S\)-模自然对应于右
\(\mathscr A(G)\)-余模.}

令 \(\mathscr A=\mathscr A(G)\). 若 \(\mathscr F\) 与
\(\mathscr F'\) 是 \(\mathscr A\)-余模, 则余模态射
\(f:\mathscr F\to\mathscr F'\) 是使下图交换的
\(\mathcal O_S\)-模态射.

% Source page transition: Exp1 component-local p. 30;
% combined-reader p. 40; re-edition printed p. 30.

\begin{figure}[H]
  \centering
  \begin{tikzcd}[column sep=4.5em,row sep=3.0em]
  \mathscr F
    \arrow[r,"f"]
    \arrow[d,"\mu_{\mathscr F}"']
  &
  \mathscr F'
    \arrow[d,"\mu_{\mathscr F'}"]
  \\
  \mathscr F\otimes\mathscr A
    \arrow[r,"f\otimes\operatorname{Id}"]
  &
  \mathscr F'\otimes\mathscr A
  \qquad .
\end{tikzcd}

  \phantomsection\label{fig:sga3-exp1-p030-comodule-morphism-square}
\end{figure}

由此得到范畴 \((\mathscr A\text{-}\mathbf{Comod.})\);
我们用 \((\mathscr A\text{-}\mathbf{Comod.q.c.})\) 表示其由
作为 \(\mathcal O_S\)-模是拟凝聚的那些 \(\mathscr A\)-余模
组成的全子范畴. 因而我们已证明如下陈述.

\medskip
\noindent\SGAInlineRefTarget{sga3:i:proposition:4-7-2}\textbf{命题 4.7.2.}
\textit{设 \(G\) 是在 \(S\) 上仿射的 \(S\)-群. 有以下范畴等价
\[
  (G\text{-}\mathcal O_S\text{-}\mathbf{Mod.})
    \simeq(\mathscr A(G)\text{-}\mathbf{Comod.}),
\]
\[
  (G\text{-}\mathcal O_S\text{-}\mathbf{Mod.q.c.})
    \simeq(\mathscr A(G)\text{-}\mathbf{Comod.q.c.}).
\]}

此外, 若 \(S=\operatorname{Spec}(\Lambda)\) 仿射, 并令
\(\Lambda[G]=\Gamma(S,\mathscr A(G))\), 则有一个范畴等价
\[
  (\mathscr A(G)\text{-}\mathbf{Comod.q.c.})
    \simeq(\Lambda[G]\text{-}\mathbf{Comod.}).
\]
\footnote{编者注: 上面一句为新增内容.}

进一步假设 \(\mathscr A=\mathscr A(G)\) 是平坦 \(\mathcal O_S\)-模.
\footnote{编者注: 下面一段由 [Se68,
\SGARefLink{sga3:i:subsection:1-3}{第 1.3 节}] 补入.}
设 \(\mathscr E\) 是一个 \(\mathscr A\)-余模, \(\mathscr F\) 是
\(\mathscr E\) 的 \(\mathcal O_S\)-子模. 由于 \(\mathscr A\)
在 \(\mathcal O_S\) 上平坦, 可以将 \(\mathscr F\otimes\mathscr A\)
(分别地, \(\mathscr F\otimes\mathscr A\otimes\mathscr A\))
等同于 \(\mathscr E\otimes\mathscr A\)
(分别地, \(\mathscr E\otimes\mathscr A\otimes\mathscr A\)) 的
\(\mathcal O_S\)-子模. 假设 \(\mu_{\mathscr E}\) 将
\(\mathscr F\) 映到 \(\mathscr F\otimes\mathscr A\) 中.
则其限制
\[
  \mu_{\mathscr F}:\mathscr F\longrightarrow
    \mathscr F\otimes\mathscr A
\]
赋予 \(\mathscr F\) 一个 \(\mathscr A\)-余模结构; 称
\(\mathscr F\) 为 \(\mathscr E\) 的\emph{子余模}.
取商后, \(\mu_{\mathscr E}\) 定义了一个 \(\mathcal O_S\)-模态射
\[
  \mathscr E/\mathscr F\longrightarrow
  (\mathscr E/\mathscr F)\otimes\mathscr A,
\]
从而赋予 \(\mathscr E/\mathscr F\) 以 \(\mathscr A\)-余模结构.
若 \(f:\mathscr E\to\mathscr E'\) 是 \(\mathscr A\)-余模态射,
则 \(\operatorname{Ker}f\) (分别地, \(\operatorname{Im}f\)) 是
\(\mathscr E\) (分别地, \(\mathscr E'\)) 的 \(\mathscr A\)-子余模,
并且 \(f\) 诱导一个 \(\mathscr A\)-余模同构
\[
  \mathscr E/\operatorname{Ker}f
    \xrightarrow{\ \sim\ }
  \operatorname{Im}f.
\]
此外, 若 \(\mathscr E\) 和 \(\mathscr E'\) 是拟凝聚
\(\mathcal O_S\)-模, 则 \(\operatorname{Ker}f\) 和
\(\operatorname{Im}f\) 也都是拟凝聚的. 因此,
\((\mathscr A\text{-}\mathbf{Comod.})\) 与
\((\mathscr A\text{-}\mathbf{Comod.q.c.})\) 都是阿贝尔范畴.

\medskip
\noindent\SGAInlineRefTarget{sga3:i:corollary:4-7-2-1}\textbf{推论 4.7.2.1.}
\textit{假设 \(G\) 在 \(S\) 上仿射且平坦. 那么, 范畴
\((G\text{-}\mathcal O_S\text{-}\mathbf{Mod.q.c.})\)
(分别地, \((G\text{-}\mathcal O_S\text{-}\mathbf{Mod.})\))
等价于在 \(\mathcal O_S\) 上拟凝聚的 \(\mathscr A(G)\)-余模范畴
(分别地, 所有 \(\mathscr A(G)\)-余模的范畴), 因而是阿贝尔范畴.}

\medskip
\noindent\SGAInlineRefTarget{sga3:i:paragraph:4-7-3}\textbf{4.7.3.}
现假设 \(G\) 是\emph{可对角化群}, 即 \(\mathscr A(G)\)
是交换群 \(M\) 在环层 \(\mathcal O_S\) 上的群代数.
若 \(\mathscr F\) 是 \(\mathcal O_S\)-模, 则
\[
  \mathscr F\otimes\mathscr A(G)
    =\coprod_{m\in M}\mathscr F\otimes m\mathcal O_S.
\]
因此, 给定一个 \(\mathcal O_S\)-模态射
\[
  \mu:\mathscr F\longrightarrow\mathscr F\otimes\mathscr A(G)
\]

% Source page transition: Exp1 component-local p. 31;
% combined-reader p. 41; re-edition printed p. 31. Only the continuation
% of Section 4.7 is included; Section 5 below it is deliberately excluded.
等价于给定 \(\mathscr F\) 的 \(\mathcal O_S\)-自同态族
\((\mu_m)_{m\in M}\), 使得对 \(\mathscr F\) 在 \(S\) 的一个开子集上
的每个截面 \(x\), 族 \((\mu_m(x))\) 都是直和
\(\coprod_{m\in M}\mathscr F\) 的截面. 这意味着, 在每个充分小的
开子集上, \(\mu_m(x)\) 的限制只有有限多个非零.

对于由
\[
  \mu(x)=\sum_{m\in M}\mu_m(x)\otimes m,
\]
定义的 \(\mu\), 条件 (CM 1), 分别地 (CM 2), 成立当且仅当
\[
  \mu_m\circ\mu_{m'}=\delta_{mm'}\mu_m,
  \qquad
  \left(\text{分别地, }\sum_{m\in M}\mu_m
    =\operatorname{Id}_{\mathscr F}\right).
\]
因此 \(\mu_m\) 是两两正交的投影算子, 且其和为恒等映射.
我们已经证明下述结果.

\medskip
\noindent\SGAInlineRefTarget{sga3:i:proposition:4-7-3}\textbf{命题 4.7.3.}
\textit{若 \(G=D_S(M)\) 是可对角化 \(S\)-群, 则拟凝聚
\(G\)-\(\mathcal O_S\)-模范畴 (分别地, 所有
\(G\)-\(\mathcal O_S\)-模的范畴) 等价于按 \(M\) 分次的拟凝聚
\(\mathcal O_S\)-模范畴 (分别地, 所有这样分次的
\(\mathcal O_S\)-模的范畴).}

\medskip
\noindent\textbf{注.}
若 \(\mathscr F\) 带有在 \(G\) 的作用下保持不变的
\(\mathcal O_S\)-代数结构, 则 \(\mathscr F\) 的分次是代数分次.
更确切地说:

\medskip
\noindent\SGAInlineRefTarget{sga3:i:corollary:4-7-3-1}\textbf{推论 4.7.3.1.}
\textit{函子 \(\mathscr A\mapsto\operatorname{Spec}(\mathscr A)\)
诱导了按 \(M\) 分次的拟凝聚 \(\mathcal O_S\)-代数范畴
与如下范畴的反范畴之间的等价: 带有 \(S\)-算子群
\(G=D_S(M)\) 且在 \(S\) 上仿射的 \(S\)-概形的范畴.}

\medskip
\noindent\SGAInlineRefTarget{sga3:i:proposition:4-7-4}\textbf{命题 4.7.4.}
\textit{设 \(G\) 是可对角化 \(S\)-群. 若
\[
  0\longrightarrow\mathscr F_1\longrightarrow\mathscr F_2
    \longrightarrow\mathscr F_3\longrightarrow0
\]
是拟凝聚 \(G\)-\(\mathcal O_S\)-模的正合列, 并且作为
\(\mathcal O_S\)-模的列是分裂的, 则它作为
\(G\)-\(\mathcal O_S\)-模的列也是分裂的.}

事实上, 若 \(G=D_S(M)\), 则每个 \(\mathscr F_i\) 都由
\((\mathscr F_i)_m\) 给出分次, 且对每个 \(m\in M\),
\(\mathcal O_S\)-模的列
\[
  0\longrightarrow(\mathscr F_1)_m\longrightarrow(\mathscr F_2)_m
    \longrightarrow(\mathscr F_3)_m\longrightarrow0
\]
都是分裂的. 由前一命题即得结论.

% End of assigned component. Section 5 (Group cohomology) begins next on
% Exp1 component-local p. 31 / combined-reader p. 41.
% Translation statistics: source_lines=829; output_lines=795; inline_math=352; display_math=58; align_environments=5.
% Review: complete assigned component; pending_TODO=0; standalone_compile=blocked_missing_ctex.sty; integrated_compile=not_run.
